\chapter{Fazit und Ausblick}
\label{ch:fazit}

% ============================================================
\section{Fazit}
% ============================================================

Diese Projektarbeit dokumentiert den vollständigen Entwicklungs-,
Integrations- und Evaluationsprozess eines Follow-the-Gap-basierten
Hindernisvermeidungssystems für das MXCarkit.

Das gesetzte Ziel -- ein Fahrzeug, das selbstständig Hindernisse
erkennt, ausweicht und bei unausweichlichen Hindernissen anhält --
konnte im Rahmen der durchgeführten Experimente \emph{nur
teilweise} erreicht werden. Die Einzelkomponenten des Stacks
(LiDAR-Empfang, Hinderniskonvertierung, Planerkern, Geschwindigkeitsadapter
und Ackermann-Ausgabe) funktionierten im Einzeltest korrekt.
Im Zusammenspiel verhinderten jedoch systemische Fehler
die kollisionsfreie Navigation.

\textbf{Zu den wesentlichen Erkenntnissen dieser Arbeit zählen:}

\begin{itemize}
    \item \textbf{Kalibrierung ist entscheidend:} Der Parameter
    \texttt{front\_center\_deg} hatte enormen Einfluss auf das
    Gesamtverhalten. Eine Abweichung von nur $45°$ (von $135°$ auf $90°$)
    führte dazu, dass der Algorithmus konstant am Grenzwert operierte
    und keinerlei situationsabhängige Reaktion zeigte.

    \item \textbf{TF-Infrastruktur ist nicht optional:}
    Reaktive Planungsalgorithmen hängen fundamental von der korrekten
    räumlichen Zuordnung der Sensordaten ab. Das Fehlen der
    \texttt{base\_link}$\to$\texttt{laser}-Transformation verhinderte
    eine automatische, robuste Kalibrierung und war die tiefste
    Ursache des Systemversagens.

    \item \textbf{Hindernismodellierung beeinflusst Lückenbewertung:}
    Ein Sicherheitsradius von $0{,}01$\,m ist für ein Fahrzeug
    der Breite $>30$\,cm unzureichend. Die Wahl des Hindernisradius
    direkt beeinflusst, welche Lücken als passierbar gelten.

    \item \textbf{Modularität erleichtert Debugging:}
    Die konsequent modulare Paketstruktur mit dedizierten
    Debug-Topics (\texttt{/visualize\_obstacles},
    \texttt{/autonomous/ftg/planner\_status}) ermöglichte eine
    systematische Eingrenzung der Fehlerquellen -- ohne diese
    Transparenz wäre das Debugging erheblich schwieriger gewesen.

    \item \textbf{Container-Grenzen sind ROS-Grenzen:}
    Das DDS-basierte Kommunikationsmodell von ROS~2 funktioniert
    gut innerhalb eines Netzwerks, jedoch entstehen durch
    unterschiedliche ROS~2-Versionen und QoS-Inkompatibilitäten
    unsichtbare Kommunikationsgrenzen.
\end{itemize}

% ============================================================
\section{Ausblick}
% ============================================================

Mit den in Kapitel~\ref{ch:diskussion} beschriebenen
Korrekturen ist zu erwarten, dass das System die ursprüngliche
Zielsetzung -- kollisionsfreies Fahren mit Nothalt --
vollständig erfüllen kann. Über die unmittelbaren Korrekturen
hinaus bieten sich folgende Erweiterungen an:

\begin{enumerate}
    \item \textbf{Clustering-Integration:}
    Ein DBSCAN-basierter Clusteringschritt zwischen
    \texttt{obstacle\_substitution} und dem Planerkern
    würde die Hindernisrepräsentation erheblich verbessern
    und das beobachtete Zickzack-Lenkverhalten reduzieren.

    \item \textbf{Systematische Parametrierung via Rosbag:}
    Die aufgezeichneten Rosbags können für offline
    Parametertests genutzt werden, ohne das reale
    Fahrzeug jedes Mal neu starten zu müssen.

    \item \textbf{Foxglove-Dashboard:}
    Eine Foxglove-Konfiguration mit 3D-Panel
    (Hindernisvisualisierung, Heading-Pfeil),
    Plot-Panel (\texttt{/final\_heading\_angle},
    \texttt{/gap\_found}) und Raw-Messages-Panel
    würde zukünftige Kalibrierung und Diagnose
    erheblich erleichtern.

    \item \textbf{Streckentest:}
    Erprobung des Systems auf einer definierten
    Teststrecke mit bekannten Hindernishindernissen
    in verschiedenen Konfigurationen (Engstellen,
    L-Kurven, dynamische Hindernisse).

    \item \textbf{Geschwindigkeitsskalierung:}
    Nach erfolgreicher Kollisionsfreiheit kann
    die \texttt{cruise\_speed} schrittweise erhöht
    und die Geschwindigkeitsstufen feinjustiert werden.
\end{enumerate}
